\documentclass[10pt]{article}
\usepackage[margin=0.75in]{geometry}
\usepackage{amsmath, amssymb, amsthm}
\usepackage{microtype}
\usepackage{enumitem}
\usepackage{fancyhdr}
\usepackage{titlesec}
\usepackage{tcolorbox}
\usepackage{booktabs}

\linespread{1.08}
\setlength{\parskip}{0.4ex plus 0.2ex minus 0.1ex}

\titleformat{\section}{\large\bfseries}{\thesection.}{0.5em}{}
\titleformat{\subsection}{\normalsize\bfseries}{\thesubsection}{0.5em}{}
\titlespacing*{\section}{0pt}{1.5ex}{1ex}
\titlespacing*{\subsection}{0pt}{1ex}{0.5ex}

\setlist{itemsep=1pt, topsep=3pt, parsep=1pt, leftmargin=1.5em}

\pagestyle{fancy}
\fancyhf{}
\fancyhead[L]{\small EECS 127/227AT Discussion 1}
\fancyhead[R]{\small Spring 2026}
\fancyfoot[C]{\small\thepage}
\renewcommand{\headrulewidth}{0.4pt}

\newcommand{\RR}{\mathbb{R}}
\newcommand{\NN}{\mathcal{N}}
\newcommand{\Range}{\mathcal{R}}
\DeclareMathOperator{\rank}{rank}
\DeclareMathOperator{\trace}{tr}

\begin{document}

\begin{center}
{\Large\bfseries EECS 127/227AT \quad Optimization Models in Engineering \quad UC Berkeley \quad Spring 2026}\\[6pt]
{\LARGE\bfseries Discussion 1}
\end{center}

\section{Invertibility of $A^\top A$}

In this problem, we show that if the matrix $A \in \RR^{m \times n}$ has a full column rank, then the matrix $A^\top A$ is invertible.

\begin{enumerate}[label=(\alph*)]
\item Show that if a vector $\vec{x}$ is in the null space of $A$ then $\vec{x}$ is in the null space of $A^\top A$.

\vspace{8cm}

\item Conversely, show that if $\vec{x}$ is in the null space of $A^\top A$ then $\vec{x}$ is in the null space of $A$.

\vspace{10cm}

\item Given that matrix $A$ has a full column rank, what can you say about its null space? What does this imply about the null space and invertibility of the matrix $A^\top A$?

\vspace{5cm}

\end{enumerate}

\newpage

\section{Eigenvalues}

Let $A \in \RR^{n \times n}$ have the eigendecomposition $P\Lambda P^{-1}$ where $\Lambda \in \RR^{n \times n}$ is a diagonal matrix with entries consisting of the eigenvalues $\lambda_1, \lambda_2, \ldots, \lambda_n$ and $P \in \RR^{n \times n}$ is an invertible matrix. Note that this is equivalent to stating that $A$ is diagonalizable via the transformation,
\begin{equation}
P^{-1}AP = \Lambda.
\end{equation}

\begin{enumerate}[label=(\alph*)]
\item Show that $A^m = P\Lambda^m P^{-1}$, for integer $m \geq 1$.

\vspace{10cm}

\item Show that determinant of $A$ is the product of its eigenvalues, i.e.
\begin{equation}
\det(A) = \prod_{i=1}^{n} \lambda_i.
\end{equation}

\textit{HINT: We have the identity} $\det(XY) = \det(X)\det(Y)$.

\vspace{8cm}

\end{enumerate}

\vfill
\noindent\textcopyright{} UCB EECS 127/227AT, Spring 2026. All Rights Reserved. This may not be publicly shared without explicit permission.

\end{document}
